\chapter{La variabile gaussiana e la funzione \texorpdfstring{$Q$}{Q}}

Fra tutte le variabili continue ce n'è una che compare ovunque: il rumore termico di un resistore, gli errori di
misura, le altezze delle persone, il rumore che disturba un segnale radio. È la \textbf{gaussiana} (o \textbf{normale}),
con la sua densità a forma di campana. Il motivo della sua onnipresenza lo vedremo nel terzo pacchetto di slide, con il
teorema centrale del limite: quando una grandezza è la somma di tantissimi piccoli contributi indipendenti (come il
rumore termico, somma degli effetti di miliardi di elettroni), la sua densità diventa gaussiana. In questo capitolo ne
studiamo la densità, la media e la varianza, e la funzione $Q$, che serve a calcolarne le probabilità. Alla fine
useremo tutto per calcolare quanti bit sbaglia un ricevitore disturbato dal rumore.

% ============================================================
\section{La gaussiana standard}
% ============================================================

\dfn{Gaussiana standard}{
  Una variabile aleatoria $X_0$, con alfabeto $\mathcal{X} = \R$, si dice \term{gaussiana} (o \term{normale})
  \term{standard}, in breve $X_0 \sim \mathcal{N}(0, 1)$, se ha densità
  \[ f_{X_0}(x_0) = \frac{1}{\sqrt{2\pi}}\, e^{-\frac{x_0^2}{2}} , \qquad x_0 \in \R . \]
}

La densità è una campana simmetrica attorno a 0, con il massimo $\frac{1}{\sqrt{2\pi}} \approx 0{,}399$ nell'origine.
Le code scendono velocissime: l'esponente $-\frac{x^2}{2}$ cresce in valore assoluto con il quadrato di $x$, e già a
$x = 4$ la densità è circa tremila volte più bassa del massimo ($e^{-8} \approx \frac{1}{3000}$). Il fattore $\frac{1}{\sqrt{2\pi}}$ non è scelto a caso: è
l'unico che rende l'area uguale a 1. Per dimostrarlo serve uno dei calcoli più belli della matematica, che passa per il
piano e per le coordinate polari.

\begin{figure}[H]
\centering
\begin{tikzpicture}
\begin{axis}[quaderno, axis lines=left, width=11cm, height=5cm, xmin=-4.2, xmax=4.4, ymin=0, ymax=0.47,
  xtick={-3,-2,-1,0,1,2,3}, ytick={0.2,0.399}, yticklabels={$0{,}2$,$\frac{1}{\sqrt{2\pi}}$}, xlabel=$x_0$,
  ylabel={$f_{X_0}(x_0)$}]
  \addplot[mspiA, domain=-4:4, samples=121] {exp(-x^2/2)/sqrt(2*pi)};
  \draw[mspigray, dashed] (axis cs:-4.2,0.399) -- (axis cs:0,0.399);
  \draw[mspigray, dashed] (axis cs:-1,0) -- (axis cs:-1,0.242);
  \draw[mspigray, dashed] (axis cs:1,0) -- (axis cs:1,0.242);
  \node[font=\scriptsize, anchor=west, align=left] at (axis cs:1.1,0.3) {in $\pm 1$ la curva cambia\\ concavità (flessi)};
\end{axis}
\end{tikzpicture}
\caption{La densità gaussiana standard: la «curva a campana».}
\end{figure}

\begin{richiamo}{le coordinate polari}
Un punto del piano si può individuare con le coordinate cartesiane $(x, y)$ oppure con le \term{coordinate polari}
$(r, \theta)$: la distanza $r$ dall'origine e l'angolo $\theta$ che il segmento dall'origine al punto forma con l'asse
$x$. Il legame è $x = r\cos\theta$, $y = r\sin\theta$, e per Pitagora $x^2 + y^2 = r^2$.

Per calcolare un integrale doppio in coordinate polari si divide il piano in «rettangolini polari»: le zone fra due
circonferenze di raggi $r$ e $r + \mathrm{d}r$ e fra due semirette di angoli $\theta$ e $\theta + \mathrm{d}\theta$.
Un rettangolino polare è quasi un rettangolo con un lato lungo $\mathrm{d}r$ (in direzione radiale) e l'altro lungo
$r\,\mathrm{d}\theta$ (un arco di circonferenza di raggio $r$: la lunghezza di un arco è raggio per angolo, in
radianti). La sua area è quindi
\[ \mathrm{d}A = r \dd r \dd\theta , \]
con quel fattore $r$ in più: i rettangolini lontani dall'origine sono più larghi.
\begin{center}
\begin{tikzpicture}[scale=1.4]
  \draw[-{Stealth}, mspigray] (-0.3,0) -- (2.6,0) node[below, font=\small, black] {$x$};
  \draw[-{Stealth}, mspigray] (0,-0.3) -- (0,2.3) node[left, font=\small, black] {$y$};
  \fill[evid] (25:1.5) arc (25:40:1.5) -- (40:1.9) arc (40:25:1.9) -- cycle;
  \draw[mspiric] (25:1.5) arc (25:40:1.5) -- (40:1.9) arc (40:25:1.9) -- cycle;
  \draw[mspigray, dashed] (0,0) -- (25:2.3); \draw[mspigray, dashed] (0,0) -- (40:2.3);
  \draw[mspigray, dashed] (1.5,0) arc (0:60:1.5); \draw[mspigray, dashed] (1.9,0) arc (0:60:1.9);
  \draw (0.5,0) arc (0:25:0.5); \node[font=\scriptsize] at (12:0.7) {$\theta$};
  \draw (0.85,0) arc (0:40:0.85); \node[font=\scriptsize, anchor=south east] at (40:0.85) {$\theta + \mathrm{d}\theta$};
  \node[font=\scriptsize, below] at (1.5,0) {$r$}; \node[font=\scriptsize, below] at (1.95,0) {$r + \mathrm{d}r$};
  \node[font=\small, anchor=west, text width=6.6cm] at (2.9,1.1)
    {Il rettangolino polare giallo ha i lati circa $\mathrm{d}r$ e $r\,\mathrm{d}\theta$: area $r\,\mathrm{d}r\,\mathrm{d}\theta$.
     Coprendo tutto il piano, $r$ va da 0 a $\infty$ e $\theta$ da 0 a $2\pi$.};
\end{tikzpicture}
\end{center}
\end{richiamo}

\teorema[gaussintegrale]{L'integrale di Gauss}{
  \[ \int_{-\infty}^{\infty} e^{-\frac{x^2}{2}} \dd x = \sqrt{2\pi} , \qquad\text{quindi}\qquad
     \int_{-\infty}^{\infty} \frac{1}{\sqrt{2\pi}}\, e^{-\frac{x^2}{2}} \dd x = 1 . \]
}

\begin{dimostrazione}[passando al quadrato e nel piano]
\passo{Il problema} La funzione $e^{-x^2/2}$ non ha una primitiva che si possa scrivere con le funzioni elementari
(polinomi, esponenziali, logaritmi, seni\ldots): l'integrale non si può calcolare direttamente. Il trucco è calcolarne il
quadrato. Sia $I = \int_{-\infty}^{\infty} e^{-x^2/2} \dd x$.
\passo{Il quadrato è un integrale doppio} Il nome della variabile di integrazione non conta, quindi
\[ I^2 = \int_{-\infty}^{\infty} e^{-\frac{x^2}{2}} \dd x \cdot \int_{-\infty}^{\infty} e^{-\frac{y^2}{2}} \dd y =
   \iint_{\R^2} e^{-\frac{x^2 + y^2}{2}} \dd x \dd y , \]
perché il prodotto di due integrali in variabili diverse è l'integrale doppio del prodotto (Capitolo~14), ed
$e^{-x^2/2}\, e^{-y^2/2} = e^{-(x^2 + y^2)/2}$.
\passo{Coordinate polari} La funzione dipende solo da $x^2 + y^2 = r^2$, la distanza dall'origine: è una campana a
simmetria circolare. In coordinate polari:
\[ I^2 = \int_0^{2\pi} \int_0^{\infty} e^{-\frac{r^2}{2}}\, r \dd r \dd\theta . \]
\passo{Ora l'integrale si fa} Grazie al fattore $r$ la primitiva c'è: per la regola della catena
$\frac{\mathrm{d}}{\mathrm{d}r}\left( -e^{-r^2/2} \right) = r\, e^{-r^2/2}$. Quindi
\[ \int_0^{\infty} e^{-\frac{r^2}{2}}\, r \dd r = \left[ -e^{-\frac{r^2}{2}} \right]_0^{\infty} = 0 - (-1) = 1 , \qquad
   I^2 = \int_0^{2\pi} 1 \dd\theta = 2\pi . \]
\passo{Conclusione} $I > 0$ (è l'integrale di una funzione positiva), quindi $I = \sqrt{2\pi}$.
\end{dimostrazione}

\teorema[gaussmomenti]{Media e varianza della gaussiana standard}{
  Se $X_0 \sim \mathcal{N}(0, 1)$: \qquad $E[X_0] = \mu_{X_0} = 0$, \qquad $\sigma_{X_0}^2 = E[X_0^2] = 1$.
}

\begin{dimostrazione}
\passo{La media} L'integrando $x\, \frac{1}{\sqrt{2\pi}} e^{-x^2/2}$ è dispari, e le due metà sono finite:
$\int_0^{\infty} x\, e^{-x^2/2} \dd x = \left[ -e^{-x^2/2} \right]_0^{\infty} = 1$. Quindi la media è 0.
\passo{Il valore quadratico medio, per parti} Si scrive $x^2 e^{-x^2/2} = x \cdot \left( x\, e^{-x^2/2} \right)$, con
$f(x) = x$ (derivata 1) e $g'(x) = x\, e^{-x^2/2}$ (primitiva $g(x) = -e^{-x^2/2}$):
\[ E[X_0^2] = \frac{1}{\sqrt{2\pi}} \int_{-\infty}^{\infty} x^2 e^{-\frac{x^2}{2}} \dd x = \frac{1}{\sqrt{2\pi}} \left\{ \Big[
   -x\, e^{-\frac{x^2}{2}} \Big]_{-\infty}^{\infty} + \int_{-\infty}^{\infty} e^{-\frac{x^2}{2}} \dd x \right\} . \]
\passo{Si conclude} Il termine di bordo è nullo (l'esponenziale batte la potenza), e l'integrale rimasto è $\sqrt{2\pi}$:
$E[X_0^2] = \frac{\sqrt{2\pi}}{\sqrt{2\pi}} = 1$. Siccome la media è 0, la varianza coincide con il valore quadratico medio.
\end{dimostrazione}

% ============================================================
\section{La gaussiana generica}
% ============================================================

Tutte le altre gaussiane si ottengono dalla standard spostandola e allargandola.

\teorema[gaussgenerica]{Densità della gaussiana generica}{
  Sia $X_0 \sim \mathcal{N}(0, 1)$ e sia $X = \sigma_X X_0 + \mu_X$, con $\sigma_X > 0$ e $\mu_X \in \R$. Allora
  \[ f_X(x) = \frac{1}{\sqrt{2\pi\sigma_X^2}}\, e^{-\frac{(x - \mu_X)^2}{2\sigma_X^2}} , \qquad x \in \R , \]
  e $E[X] = \mu_X$, $\mathrm{VAR}(X) = \sigma_X^2$. Si scrive $X \sim \mathcal{N}(\mu_X, \sigma_X^2)$.
}

\begin{dimostrazione}
\passo{La densità} È la trasformazione lineare $g(x) = \sigma_X x + \mu_X$, crescente. Con la formula del
Capitolo~12, $f_X(x) = \frac{1}{\sigma_X} f_{X_0}\left( \frac{x - \mu_X}{\sigma_X} \right)$:
\[ f_X(x) = \frac{1}{\sigma_X}\, \frac{1}{\sqrt{2\pi}}\, e^{-\frac{1}{2}\left( \frac{x - \mu_X}{\sigma_X} \right)^2} =
   \frac{1}{\sqrt{2\pi\sigma_X^2}}\, e^{-\frac{(x - \mu_X)^2}{2\sigma_X^2}} . \]
\passo{La media} Per la linearità: $E[X] = E[\sigma_X X_0 + \mu_X] = \sigma_X E[X_0] + \mu_X = 0 \cdot \sigma_X + \mu_X =
\mu_X$.
\passo{La varianza}
\[ \mathrm{VAR}(X) = E[(X - \mu_X)^2] = E[(\sigma_X X_0 + \mu_X - \mu_X)^2] = E[(\sigma_X X_0)^2] = \sigma_X^2\, E[X_0^2] =
   \sigma_X^2 . \]
\end{dimostrazione}

I due parametri hanno quindi un significato diretto: $\mu_X$ è il centro della campana (e la sua media), $\sigma_X$ la
sua larghezza (e la sua deviazione standard). Una gaussiana è completamente determinata dalla coppia
$(\mu_X, \sigma_X^2)$. Il procedimento si può anche rovesciare: se $X \sim \mathcal{N}(\mu_X, \sigma_X^2)$, la variabile
\textbf{standardizzata} $X_0 = \frac{X - \mu_X}{\sigma_X}$ è una gaussiana standard.

\begin{figure}[H]
\centering
\begin{tikzpicture}
\begin{axis}[quaderno, axis lines=left, width=12cm, height=5.2cm, xmin=-4.4, xmax=6.4, ymin=0, ymax=0.85,
  xtick={-3,-2,-1,0,1,2,3,4,5}, ytick={0.4,0.8}, yticklabels={$0{,}4$,$0{,}8$}, xlabel=$x$, ylabel={$f_X(x)$},
  legend style={font=\small, at={(0.99,0.98)}, anchor=north east}]
  \addplot[mspiA, domain=-4.2:6.2, samples=161] {exp(-x^2/2)/sqrt(2*pi)};
  \addplot[mspiB, domain=-4.2:6.2, samples=161] {exp(-x^2/(2*0.25))/sqrt(2*pi*0.25)};
  \addplot[mspiC, domain=-4.2:6.2, samples=161] {exp(-(x-2)^2/(2*2.25))/sqrt(2*pi*2.25)};
  \legend{{$\mathcal{N}(0, 1)$}, {$\mathcal{N}(0;\ 0{,}25)$}, {$\mathcal{N}(2;\ 2{,}25)$}}
\end{axis}
\end{tikzpicture}
\caption{Tre gaussiane. Con $\sigma = 0{,}5$ la campana è stretta e alta, con $\sigma = 1{,}5$ larga e bassa: l'area resta
sempre 1. La media $\mu$ sposta il centro.}
\end{figure}

% ============================================================
\section{La funzione \texorpdfstring{$Q$}{Q}}
% ============================================================

Per calcolare la probabilità che una gaussiana cada in un intervallo servirebbe la sua CDF. Ma né la CDF né la CCDF di
$X_0 \sim \mathcal{N}(0, 1)$ sono note in forma esplicita, perché $e^{-\gamma x^2}$ non ammette primitive elementari
(è lo stesso ostacolo incontrato con l'integrale di Gauss, ma ora l'integrale non è su tutta la retta, e il trucco del
quadrato non funziona). Si dà allora un nome alla funzione che serve, e se ne tabulano i valori.

\dfn{Funzione Q}{
  La \term{funzione $Q$} è la CCDF della gaussiana standard, cioè la probabilità che $X_0 \sim \mathcal{N}(0, 1)$
  superi il valore $x$:
  \[ Q(x) \stackrel{\triangle}{=} P(X_0 > x) = 1 - F_{X_0}(x) = \frac{1}{\sqrt{2\pi}} \int_x^{+\infty} e^{-\frac{t^2}{2}} \dd t . \]
}

\begin{figure}[H]
\centering
\begin{tikzpicture}
\begin{axis}[quaderno, axis y line=none, width=10cm, height=4.8cm, xmin=-4.2, xmax=4.4, ymin=0, ymax=0.47,
  xtick={-3,-2,-1,0,1.3,2,3}, xticklabels={$-3$,$-2$,$-1$,$0$,$x$,$2$,$3$}, ytick=\empty, xlabel=$t$]
  \addplot[fill=evid, draw=none, domain=1.3:4, samples=40] {exp(-x^2/2)/sqrt(2*pi)} \closedcycle;
  \addplot[mspiA, domain=-4:4, samples=121] {exp(-x^2/2)/sqrt(2*pi)};
  \node[font=\small, align=center] at (axis cs:-0.25,0.13) {$F_{X_0}(x)$\\ $= 1 - Q(x)$};
  \draw[-{Stealth}] (axis cs:2.6,0.2) node[above, font=\small] {$Q(x)$} -- (axis cs:1.75,0.05);
\end{axis}
\end{tikzpicture}
\caption{$Q(x)$ è l'area della coda destra della gaussiana standard, oltre $x$.}
\end{figure}

Con la funzione $Q$ si calcola qualsiasi probabilità gaussiana. Per la standard, $F_{X_0}(x) = 1 - Q(x)$ e la probabilità
di un intervallino è
\[ P_{X_0}(x; \Delta x) = Q\left( x - \frac{\Delta x}{2} \right) - Q\left( x + \frac{\Delta x}{2} \right) . \]
Per una gaussiana qualsiasi si standardizza.

\teorema[qgenerica]{Le code di una gaussiana qualsiasi}{
  Se $X \sim \mathcal{N}(\mu_X, \sigma_X^2)$:
  \[ P(X > x) = 1 - F_X(x) = Q\left( \frac{x - \mu_X}{\sigma_X} \right) . \]
}

\begin{dimostrazione}
\passo{} Si scrive $X = X_0 \sigma_X + \mu_X$ e si isola $X_0$ (dividere per $\sigma_X > 0$ non cambia il verso della
disuguaglianza):
\[ P(X > x) = P(X_0 \sigma_X + \mu_X > x) = P\left( X_0 > \frac{x - \mu_X}{\sigma_X} \right) =
   \int_{\frac{x - \mu_X}{\sigma_X}}^{+\infty} \frac{1}{\sqrt{2\pi}}\, e^{-\frac{t^2}{2}} \dd t = Q\left( \frac{x -
   \mu_X}{\sigma_X} \right) . \]
\end{dimostrazione}

La quantità $\frac{x - \mu_X}{\sigma_X}$ dice \textbf{di quante deviazioni standard} il valore $x$ si trova sopra la media:
è tutto quello che conta.

\subsection{Le proprietà della funzione \texorpdfstring{$Q$}{Q}}

\teorema[qproprieta]{Proprietà della funzione Q}{
  \begin{enumerate}[label=\alph*)]
    \item $Q(-\infty) = 1$, \quad $Q(0) = \frac12$, \quad $Q(+\infty) = 0$;
    \item $\dfrac{\dd Q(x)}{\dd x} = -\dfrac{1}{\sqrt{2\pi}}\, e^{-\frac{x^2}{2}} < 0$ per ogni $x$: $Q$ è decrescente;
    \item simmetria: $Q(-x) = 1 - Q(x)$;
    \item se $X \sim \mathcal{N}(\mu_X, \sigma_X^2)$: $P(X \geq \eta) = Q\left( \frac{\eta - \mu_X}{\sigma_X} \right)$.
  \end{enumerate}
}

\begin{dimostrazione}
\passo{a)} $Q(-\infty)$ è l'area sotto tutta la campana, cioè 1; $Q(0)$ è metà dell'area, per la simmetria; $Q(+\infty)$
è l'area di una coda vuota.
\passo{b)} Si scrive $Q(x) = 1 - \int_{-\infty}^{x} \frac{1}{\sqrt{2\pi}} e^{-t^2/2} \dd t$ e si deriva con il teorema
fondamentale del calcolo integrale: la derivata della funzione integrale è l'integrando calcolato in $x$, con il segno
meno. Spostando $x$ a destra, la coda si accorcia.
\passo{c)} Per la parità della funzione integranda, la coda sinistra oltre $-x$ è uguale alla coda destra oltre $x$:
$\int_{-\infty}^{-x} \frac{1}{\sqrt{2\pi}} e^{-t^2/2} \dd t = \int_{x}^{\infty} \frac{1}{\sqrt{2\pi}} e^{-t^2/2} \dd t = Q(x)$.
Quindi
\[ Q(-x) = \int_{-x}^{+\infty} \frac{1}{\sqrt{2\pi}} e^{-\frac{t^2}{2}} \dd t = 1 - \underbrace{\int_{-\infty}^{-x}
   \frac{1}{\sqrt{2\pi}} e^{-\frac{t^2}{2}} \dd t}_{Q(x)} = 1 - Q(x) . \]
\passo{d)} È il teorema precedente: per una variabile continua $P(X \geq \eta) = P(X > \eta)$.
\end{dimostrazione}

Grazie alla simmetria basta tabulare $Q(x)$ per $x \geq 0$. Ecco alcuni valori:
\begin{center}
\renewcommand{\arraystretch}{1.25}
\begin{tabular}{@{}c|cccccccc@{}}
  $x$ & 0 & 0,5 & 1 & 1,5 & 2 & 2,5 & 3 & 4 \\ \hline
  $Q(x)$ & 0,5 & 0,309 & 0,159 & 0,0668 & 0,0228 & 0,00621 & 0,00135 & $3{,}17 \cdot 10^{-5}$
\end{tabular}
\end{center}

\teorema[qmaggiorazione]{Una maggiorazione della funzione Q}{
  Per ogni $x > 0$:
  \[ Q(x) < \frac{1}{x\sqrt{2\pi}}\, e^{-\frac{x^2}{2}} , \]
  e per $x$ grande il secondo membro è anche una buona approssimazione: $Q(x) \simeq \frac{1}{x\sqrt{2\pi}}
  e^{-\frac{x^2}{2}}$. In particolare, per $x > \frac{1}{\sqrt{2\pi}} \approx 0{,}4$, $Q(x) < e^{-\frac{x^2}{2}}$.
}

\begin{dimostrazione}[lo stesso trucco della disuguaglianza di Markov]
\passo{Si moltiplica per un numero maggiore di 1} Nella coda $t > x > 0$, quindi $\frac{t}{x} > 1$:
\[ Q(x) = \frac{1}{\sqrt{2\pi}} \int_x^{\infty} e^{-\frac{t^2}{2}} \dd t < \frac{1}{\sqrt{2\pi}} \int_x^{\infty}
   \frac{t}{x}\, e^{-\frac{t^2}{2}} \dd t . \]
\passo{Ora la primitiva c'è} Come nell'integrale di Gauss, $t\, e^{-t^2/2}$ ha primitiva $-e^{-t^2/2}$:
\[ \frac{1}{x\sqrt{2\pi}} \int_x^{\infty} t\, e^{-\frac{t^2}{2}} \dd t = \frac{1}{x\sqrt{2\pi}} \left[ -e^{-\frac{t^2}{2}}
   \right]_x^{\infty} = \frac{1}{x\sqrt{2\pi}}\, e^{-\frac{x^2}{2}} . \]
\passo{Perché per $x$ grande è anche un'approssimazione} Il fattore $\frac{t}{x}$ è vicino a 1 nella parte della coda
che conta, appena oltre $x$; più in là la densità è trascurabile. Quindi la maggiorazione esagera di poco. Infine, se
$x\sqrt{2\pi} > 1$, il fattore $\frac{1}{x\sqrt{2\pi}}$ è minore di 1.
\end{dimostrazione}

\begin{figure}[H]
\centering
\begin{tikzpicture}
\begin{axis}[width=11cm, height=5.6cm, ymode=log, xmin=0, xmax=5, ymin=1e-7, ymax=1.5, xtick={0,1,2,3,4,5},
  xlabel=$x$, tick label style={font=\scriptsize}, label style={font=\small}, axis lines=left,
  grid=major, grid style={mspigray!20}, legend style={font=\scriptsize, at={(0.02,0.03)}, anchor=south west}]
  \addplot[mspiA, line width=1.2pt] coordinates {(0.00,5.0000e-01) (0.25,4.0129e-01) (0.50,3.0854e-01) (0.75,2.2663e-01) (1.00,1.5866e-01) (1.25,1.0565e-01) (1.50,6.6807e-02) (1.75,4.0059e-02) (2.00,2.2750e-02) (2.25,1.2224e-02) (2.50,6.2097e-03) (2.75,2.9798e-03) (3.00,1.3499e-03) (3.25,5.7703e-04) (3.50,2.3263e-04) (3.75,8.8417e-05) (4.00,3.1671e-05) (4.25,1.0689e-05) (4.50,3.3977e-06) (4.75,1.0171e-06) (5.00,2.8665e-07)};
  \addplot[mspiB, dashed, line width=1pt, domain=0.3:5, samples=60] {exp(-x^2/2)/(x*sqrt(2*pi))};
  \addplot[mspiC, dotted, line width=1.2pt, domain=0:5, samples=60] {exp(-x^2/2)};
  \legend{$Q(x)$, $\frac{1}{x\sqrt{2\pi}} e^{-x^2/2}$, $e^{-x^2/2}$}
\end{axis}
\end{tikzpicture}
\caption{La funzione $Q$ in scala logaritmica (ogni tacca verticale è un fattore 10) e le due maggiorazioni. Da $x
\approx 2$ in poi la prima è quasi indistinguibile da $Q(x)$; la seconda è più grossolana ma più semplice.}
\end{figure}

\subsection{La regola delle tre sigma}

Quanto spesso una gaussiana si allontana dalla sua media di più di $k$ deviazioni standard? Standardizzando e usando la
simmetria (le due code sono uguali):
\[ P(\abs{X - \mu_X} > k\,\sigma_X) = P(\abs{X_0} > k) = 2\,Q(k) , \qquad P(\abs{X - \mu_X} \leq k\,\sigma_X) = 1 - 2\,Q(k) . \]

\begin{figure}[H]
\centering
\begin{tikzpicture}
\begin{axis}[quaderno, axis y line=none, width=12cm, height=5.6cm, xmin=-4.2, xmax=4.4, ymin=0, ymax=0.52,
  xtick={-3,-2,-1,0,1,2,3}, xticklabels={$\mu-3\sigma$,$\mu-2\sigma$,$\mu-\sigma$,$\mu$,$\mu+\sigma$,$\mu+2\sigma$,$\mu+3\sigma$},
  ytick=\empty, xticklabel style={font=\tiny}]
  \addplot[fill=mspiA!12, draw=none, domain=-3:3, samples=60] {exp(-x^2/2)/sqrt(2*pi)} \closedcycle;
  \addplot[fill=mspiA!25, draw=none, domain=-2:2, samples=60] {exp(-x^2/2)/sqrt(2*pi)} \closedcycle;
  \addplot[fill=mspiA!45, draw=none, domain=-1:1, samples=60] {exp(-x^2/2)/sqrt(2*pi)} \closedcycle;
  \addplot[mspiA, domain=-4:4, samples=121] {exp(-x^2/2)/sqrt(2*pi)};
  \node[font=\scriptsize, align=center] at (axis cs:0,0.15) {68,3\%};
  \draw[mspigray, densely dotted, line width=0.5pt] (axis cs:-2,0) -- (axis cs:-2,0.43) (axis cs:2,0) -- (axis cs:2,0.43)
    (axis cs:-3,0) -- (axis cs:-3,0.5) (axis cs:3,0) -- (axis cs:3,0.5);
  \draw[{Stealth}-{Stealth}, mspigray, line width=0.6pt] (axis cs:-2,0.43) -- (axis cs:2,0.43)
    node[pos=0.75, above, font=\scriptsize, text=black] {95,4\%};
  \draw[{Stealth}-{Stealth}, mspigray, line width=0.6pt] (axis cs:-3,0.5) -- (axis cs:3,0.5)
    node[pos=0.8, above, font=\scriptsize, text=black] {99,7\%};
\end{axis}
\end{tikzpicture}
\caption{La regola 68--95--99,7: le probabilità che una gaussiana cada entro una, due e tre deviazioni standard dalla
media.}
\end{figure}

\begin{center}
\renewcommand{\arraystretch}{1.25}
\begin{tabular}{@{}c|ccc@{}}
  $k$ & 1 & 2 & 3 \\ \hline
  $P(\abs{X - \mu_X} \leq k\sigma_X)$ per la gaussiana: $1 - 2Q(k)$ & 0,683 & 0,954 & 0,997 \\
  minimo garantito per qualsiasi densità (Chebyshev): $1 - \frac{1}{k^2}$ & 0 & 0,75 & 0,889
\end{tabular}
\end{center}

Una gaussiana esce dalla fascia $\mu \pm 3\sigma$ solo 3 volte su mille: in pratica, un valore a più di tre deviazioni
standard dalla media è un evento eccezionale. Ma attenzione: vale per la gaussiana, non per ogni variabile (per una
densità qualsiasi, Chebyshev garantisce solo l'89\%).

% ============================================================
\section{Un'applicazione: bit sbagliati per colpa del rumore}
% ============================================================

\begin{esempio}{La probabilità di errore di una trasmissione binaria}
Un trasmettitore invia un bit codificandolo come una tensione: $+A$ per l'1, $-A$ per lo 0, con i due bit
equiprobabili. Il canale aggiunge un rumore gaussiano $N \sim \mathcal{N}(0, \sigma^2)$, indipendente dal bit, e il
ricevitore osserva $Y = X + N$. Decide «1» se $Y > 0$ e «0» se $Y \leq 0$. Con che probabilità sbaglia?
\begin{itemize}
  \item Se è stato trasmesso l'1, $Y = A + N$: per la trasformazione lineare $Y \sim \mathcal{N}(A, \sigma^2)$. Il
        ricevitore sbaglia se $Y \leq 0$:
        \[ P(\text{errore} \mid X = A) = P(A + N \leq 0) = P(N \leq -A) = P(N \geq A) = Q\left( \frac{A}{\sigma} \right) , \]
        dove si è usata la simmetria del rumore (le code sono uguali) e la proprietà d) della funzione $Q$.
  \item Se è stato trasmesso lo 0, per simmetria l'errore ($Y > 0$) ha la stessa probabilità.
  \item Per la probabilità totale: $P_e = \frac12 Q\left(\frac{A}{\sigma}\right) + \frac12 Q\left(\frac{A}{\sigma}\right) =
        Q\left(\frac{A}{\sigma}\right)$.
\end{itemize}
\begin{center}
\begin{tikzpicture}
\begin{axis}[quaderno, axis y line=none, width=11cm, height=4.8cm, xmin=-4.6, xmax=4.8, ymin=0, ymax=0.47,
  xtick={-1.5,0,1.5}, xticklabels={$-A$,$0$,$A$}, ytick=\empty, xlabel=$y$, legend style={font=\small,
  at={(0.99,0.98)}, anchor=north east}]
  \addplot[fill=mspiB!35, draw=none, domain=-4.4:0, samples=50, forget plot] {exp(-(x-1.5)^2/2)/sqrt(2*pi)} \closedcycle;
  \addplot[fill=mspiA!35, draw=none, domain=0:4.4, samples=50, forget plot] {exp(-(x+1.5)^2/2)/sqrt(2*pi)} \closedcycle;
  \addplot[mspiA, domain=-4.4:4.4, samples=121] {exp(-(x+1.5)^2/2)/sqrt(2*pi)};
  \addplot[mspiB, domain=-4.4:4.4, samples=121] {exp(-(x-1.5)^2/2)/sqrt(2*pi)};
  \draw[black, dashed] (axis cs:0,0) -- (axis cs:0,0.44) node[above, font=\scriptsize] {soglia};
  \legend{$f_{Y \mid X = -A}$, $f_{Y \mid X = A}$}
\end{axis}
\end{tikzpicture}
\end{center}
Le due campane sono le densità di $Y$ sapendo quale bit è stato trasmesso; le code colorate che oltrepassano la soglia
sono gli errori. Più le campane sono lontane in rapporto alla loro larghezza, cioè più è grande $\frac{A}{\sigma}$, meno
si sbaglia:
\begin{center}
\renewcommand{\arraystretch}{1.25}
\begin{tabular}{@{}c|ccccc@{}}
  $A/\sigma$ & 1 & 2 & 3 & 4 & 5 \\ \hline
  $P_e = Q(A/\sigma)$ & 0,16 & 0,023 & 0,0013 & $3{,}2 \cdot 10^{-5}$ & $2{,}9 \cdot 10^{-7}$
\end{tabular}
\end{center}
Raddoppiare l'ampiezza da $\sigma$ a $2\sigma$ divide gli errori per 7; passare da $2\sigma$ a $4\sigma$ li divide per
700: la coda gaussiana crolla come $e^{-x^2/2}$. Il rapporto $A/\sigma$ è legato al \emph{rapporto segnale-rumore}, e
decidere nel modo migliore che cosa è stato trasmesso è il problema della teoria della decisione, che affronteremo nel
terzo pacchetto di slide.
\end{esempio}

\gbox{In sintesi}{mspiA!80!black}{
\begin{itemize}
  \item \textbf{Gaussiana standard} $\mathcal{N}(0, 1)$: $f(x) = \frac{1}{\sqrt{2\pi}} e^{-x^2/2}$; l'area 1 segue
        dall'integrale di Gauss $\int e^{-x^2/2} = \sqrt{2\pi}$ (in coordinate polari); media 0, varianza 1.
  \item \textbf{Gaussiana generica} $X = \sigma_X X_0 + \mu_X \sim \mathcal{N}(\mu_X, \sigma_X^2)$:
        $f_X(x) = \frac{1}{\sqrt{2\pi\sigma_X^2}} e^{-(x - \mu_X)^2/(2\sigma_X^2)}$, determinata da media e varianza;
        standardizzando, $\frac{X - \mu_X}{\sigma_X} \sim \mathcal{N}(0, 1)$.
  \item \textbf{Funzione $Q$}: $Q(x) = P(X_0 > x)$, la coda della standard; $P(X > x) = Q\left(\frac{x -
        \mu_X}{\sigma_X}\right)$. È decrescente, $Q(0) = \frac12$, $Q(-x) = 1 - Q(x)$, e
        $Q(x) < \frac{1}{x\sqrt{2\pi}} e^{-x^2/2}$ per $x > 0$.
  \item Regola 68--95--99,7: $P(\abs{X - \mu} \leq k\sigma) = 1 - 2Q(k)$.
  \item Trasmissione binaria $\pm A$ con rumore gaussiano: $P_e = Q(A/\sigma)$.
\end{itemize}
}
